% FIGURE NEEDED: setup sketch at the head of the official solution (2021 TQ14,
% solutions page 1): dashed circle "Orbit of venus" around the Sun, Venus of
% mass m on the circle at distance R_0 from the Sun, and the comet of mass
% "alpha m" approaching radially from the left.

\begin{description}
\item[14.1] Let $M$ be the mass of the Sun, then
  \begin{align*}
    \frac{GMm}{R_0^2} &= \frac{mv_0^2}{R_0} \\
    v_0^2 &= \frac{GM}{R_0}
  \end{align*}
  \hfill[1 point]

\item[14.2] The mechanical energy of Venus (before collision) is:
  \[
    E_i = -\frac{GMm}{R_0} + \frac{1}{2}mv_0^2 = -\frac{GMm}{2R_0}
  \]
  \hfill[1 point]

\item[14.3] Since the comet moves radially towards the Sun, it has no angular
  momentum (with respect to the origin in the Sun). Then, the angular momentum
  of Venus is the same as that of ``Venus-2''
  \[
    L = R_0 m v_0
  \]
  \hfill[2 point] % sic: "[2 point]"

  This allows us to find the angular component of the velocity of Venus,
  $v_\theta$, just after the collision. The angular momentum just after the
  collision is
  \[
    L = R_0 m v_0 = R_0(m + \alpha m)v_\theta
  \]
  \hfill[1 points] % sic: "[1 points]"

  Since the angular momentum is conserved it follows that
  \[
    v_\theta = \frac{v_0}{1 + \alpha}
  \]
  The comet has zero energy, so
  \[
    K = -U = +\frac{GM\alpha m}{R_0} = \frac{1}{2}\alpha m v_c^2
  \]
  \hfill[2 points]

  where $v_c$ is the velocity of the comet just before the collision. It
  follows that
  \[
    v_c^2 = \frac{2GM}{R_0} = 2v_0^2
  \]
  \hfill[1 point]

  The collision between the comet and Venus is inelastic, thus the total
  energy is not conserved. Also notice that the mass of ``Venus-2'' is
  $(1 + \alpha)m$. \hfill[2 point] % sic: "[2 point]"

  We now apply the conservation of linear momentum. Along the radial direction
  we get
  \[
    \alpha m v_c = (m + \alpha m)v_r
  \]
  \hfill[1 point]

  where $v_r$ is the velocity of ``Venus-2'' just after the collision. It
  follows that
  \[
    v_r = \frac{\alpha}{1 + \alpha}v_c = \frac{\sqrt{2}\,\alpha}{1 + \alpha}v_0
  \]
  \hfill[1 point]

\item[14.4] Venus-2 mechanical energy (we evaluate it just after the
  collision) is:
  \[
    E_f = U + K = -\frac{GMm(1 + \alpha)}{R_0}
      + \frac{1}{2}m(1 + \alpha)\left(v_\theta^2 + v_r^2\right)
  \]
  \hfill[4 points]
  \begin{align*}
    &= -(1 + \alpha)\frac{GMm}{R_0} + \frac{1}{2}(1 + \alpha)
      \left[2\left(\frac{\alpha}{1 + \alpha}\right)^2
      + \frac{1}{(1 + \alpha)^2}\right]\frac{GMm}{R_0} \\
    &= -\frac{GMm}{2R_0}\left(\frac{1 + 4\alpha}{1 + \alpha}\right)
      = E_i\left(\frac{1 + 4\alpha}{1 + \alpha}\right)
  \end{align*}
  \hfill[1 point]

\item[14.5] ``Venus-2'' orbit is no longer a circle. Since the energy is
  negative it must, therefore, be elliptical. It follows that
  \hfill[1 points] % sic: "[1 points]"
  \[
    E_f = -G\frac{M(1 + \alpha)m}{2a_f}
  \]
  \hfill[2 points]

  where $a_f$ is the semi-major axis of the orbit of Venus-2. Now, solving for
  $a_f$:
  \begin{align*}
    -G\frac{M(1 + \alpha)m}{2a_f}
      &= E_i\left(\frac{1 + 4\alpha}{1 + \alpha}\right)
       = -G\frac{Mm}{2R_0}\left(\frac{1 + 4\alpha}{1 + \alpha}\right) \\
    \frac{(1 + \alpha)}{a_f}
      &= \frac{1}{R_0}\left(\frac{1 + 4\alpha}{1 + \alpha}\right) \\
    a_f &= R_0\,\frac{(1 + \alpha)^2}{1 + 4\alpha}
  \end{align*}
  \hfill[2 points]

\item[14.6] Using Kepler's third law, we have that
  \[
    \frac{T_f}{T_0} = \left(\frac{a_f}{R_0}\right)^{3/2}
  \]
  \hfill[2 points]
  \[
    \frac{T_f}{T_0}
      = \left(\frac{(1 + \alpha)^2}{1 + 4\alpha}\right)^{3/2}
      = \frac{(1 + \alpha)^3}{(1 + 4\alpha)^{3/2}}
  \]
  \hfill[1 points] % sic: "[1 points]"

  For $\alpha < 2$, the year is shorter, for $\alpha > 2$, the year is longer,
  while for $\alpha = 2$ the duration of the year does not change.

\item[14.7] For the `Venus-2' to collide with the Sun, the perihelion radius
  of the post-collision orbit should be $R_\odot$. \hfill[1 point]
  \[
    r_p = R_\odot = 6.955\times10^{8}\,m
      = \frac{6.955\times10^{8}}{0.723\times1.496\times10^{11}}R_0
      = 0.00643R_0
  \]
  Let $x = \dfrac{R_\odot}{R_0} = 0.00643$.

  Now, we apply conservation of angular momentum right after the collision and
  at the moment ``Venus-2'' is at the perihelion:
  \[
    L = mv_0R_0 = (1 + \alpha)mv_pR_\odot
  \]
  where $v_p$ is the speed of Venus-2 at the perihelion. Solving for $v_p$
  \[
    v_p = \frac{v_0}{x(1 + \alpha_c)}
  \]
  Applying conservation of energy between the two same instants we obtain:
  \begin{align*}
    E_f &= -\frac{GM(1 + \alpha_c)m}{R_\odot}
      + \frac{1}{2}m(1 + \alpha_c)v_p^2 \\
    -\left(\frac{1 + 4\alpha_c}{1 + \alpha_c}\right)\frac{GMm}{2R_0}
      &= -\frac{GM(1 + \alpha_c)m}{xR_0}
      + \frac{1}{2}m(1 + \alpha_c)
        \left(\frac{v_0}{x(1 + \alpha_c)}\right)^2 \\
    -\left(\frac{1 + 4\alpha_c}{1 + \alpha_c}\right)\frac{GMm}{2R_0}
      &= -\frac{(1 + \alpha_c)}{x}\frac{GMm}{R_0}
      + \frac{1}{2}(1 + \alpha_c)
        \left(\frac{1}{x(1 + \alpha_c)}\right)^2\frac{GMm}{R_0} \\
    -\frac{1}{2}\left(\frac{1 + 4\alpha_c}{1 + \alpha_c}\right)
      &= -\frac{(1 + \alpha_c)}{x}
      + \frac{1}{2}(1 + \alpha_c)
        \left(\frac{1}{x(1 + \alpha_c)}\right)^2 \\
    (1 + 4\alpha_c) &= \frac{2(1 + \alpha_c)^2}{x} - \frac{1}{x^2}
  \end{align*}
  \hfill[3 points]
  \begin{align*}
    2x\alpha_c^2 + 4x(1 - x)\alpha_c - (1 - x)^2 &= 0 \\
    \alpha_c = (1 - x)\left(\sqrt{1 + \frac{1}{2x}} - 1\right) & \\
    \alpha_c = 7.824 &
  \end{align*}
  \hfill[1 point]

\item[14.8] For post-collision orbit,
  \begin{align*}
    v_\theta &= \frac{1}{1 + \alpha_c}v_0 \\
    v_r &= \frac{\sqrt{2}\,\alpha_c}{1 + \alpha_c}v_0 \\
    v_f &= \sqrt{v_r^2 + v_\theta^2}
  \end{align*}
  \hfill[1 point]
  \[
    v_f = \frac{\sqrt{2\alpha_c^2 + 1}}{1 + \alpha_c}v_0
  \]
  \hfill[3 point] % sic: "[3 point]"

  Thus,
  \begin{align*}
    \delta v &= v_0 - v_f \\
      &= \left(1 - \frac{\sqrt{2\alpha_c^2 + 1}}{1 + \alpha_c}\right)v_0 \\
      &= -0.259\,v_0
  \end{align*}
  \[
    \frac{\delta v}{v_0} = 25.9\%
  \]
  \hfill[0.5 points]
  \begin{align*}
    \tan\delta\theta &= \left(\frac{v_r}{v_\theta}\right) \\
    \tan\delta\theta &= \left(\sqrt{2}\,\alpha_c\right) \\
    \delta\theta &= 1.480\ \mathrm{rad}
  \end{align*}
  \hfill[0.5 points]
  \[
    \delta\theta = 84.84^\circ
  \]
\end{description}
